% ROOT V133 research formula summary. ONLY c=1; hbar and G remain dimensionful.
\documentclass[11pt]{article}
\usepackage{amsmath,amssymb,geometry}
\geometry{margin=1in}
\begin{document}
\title{ROOT V133: Conditional Helix Constraints, Photon Holonomy,\ Kaluza--Klein Geometry and One-loop Coupling Audit}
\date{2026-11-24}
\maketitle
\noindent\textbf{Status:} Mathematical and numerical consistency checks, not a proof of unified field theory. \\
\textbf{Unit convention:} $c=1$ ONLY; $\hbar$ and $G$ are dimensionful and independent. $[\hbar]=ML$, $[G_4]=L/M$.
\section*{1. Null helix conditional no-go}
\begin{align}
 x^\mu(t)&=(t,R\cos\Omega t,R\sin\Omega t,ut),\\
 -(\dot t)^2+\dot{\mathbf x}^2&=-1+R^2\Omega^2+u^2=0,\\
 R_0>0,\quad \Omega=\omega=k&\quad\Longrightarrow\quad R_0 k\le 1,\\
 \ddot x^\mu&=(0,-R\Omega^2\cos\Omega t,-R\Omega^2\sin\Omega t,0),\\
 \eta_{\mu\nu}\ddot x^\mu\ddot x^\nu&=R^2\Omega^4>0\quad(R\Omega\ne0).
\end{align}
A nondegenerate null helix is not a free Minkowski null geodesic. This does not constrain quantum helicity or field-vector rotation.
\section*{2. Berry holonomy of photon helicity}
\begin{align}
 \mathbf e_\sigma&=(\mathbf e_\theta+i\sigma\mathbf e_\phi)/\sqrt 2,\quad\sigma\in\{-1,+1\},\\
 \mathcal A_\sigma&=i\mathbf e_\sigma^*\cdot d\mathbf e_\sigma=\sigma\cos\theta\,d\phi,\\
 \mathcal F_\sigma&=-\sigma\sin\theta\,d\theta\wedge d\phi,\\
 \gamma_\sigma(C)&=-\sigma\Omega_{\rm solid}(C)\pmod{2\pi}.
\end{align}
\section*{3. Circle KK reduction}
\begin{align}
 ds_5^2&=g_{\mu\nu}dx^\mu dx^\nu+R^2(dy+\kappa A_\mu dx^\mu)^2,\\
 G_4&=\frac{G_5}{2\pi R},\quad p_5=\frac{n\hbar}{R},\quad [G_5]=L^2/M.
\end{align}
$G_4$ is not derived without independent inputs for $G_5$ and $R$.
\section*{4. Minimal SM one-loop gauge running}
\begin{align}
 (b_1,b_2,b_3)&=(41/10,-19/6,-7),\\
 \alpha_i^{-1}(\mu)&=\alpha_i^{-1}(M_Z)-\frac{b_i}{2\pi}\ln\frac{\mu}{M_Z},\\
 \alpha_1^{-1}(M_Z)&=\tfrac 35(1-\sin^2\theta_W)\alpha_{em}^{-1}(M_Z),\\
 \alpha_2^{-1}(M_Z)&=\sin^2\theta_W\alpha_{em}^{-1}(M_Z),\\
 \alpha_3^{-1}(M_Z)&=1/\alpha_s(M_Z).
\end{align}
At reference inputs $(M_Z,\alpha_{em}^{-1},\sin^2\theta_W,\alpha_s)=(91.1876\,{\rm GeV},127.955,0.23122,0.1180)$,
pairwise intersections are near $1.032\cdot10^{13}$, $2.433\cdot10^{14}$, and $9.725\cdot10^{16}$ GeV. No single intersection at this order. These are illustrative scheme-approximate inputs, not a full precision fit.
\section*{5. Anomaly and dimensional gatekeeping}
\begin{align}
 F_{\mu\nu}^{\rm naive}&=\nabla_\mu A_\nu-\nabla_\nu A_\mu
 =\partial_\mu A_\nu-\partial_\nu A_\mu-T^\rho{}_{\mu\nu} A_\rho,\\
 \delta_\chi F_{\mu\nu}^{\rm naive}&=-T^\rho{}_{\mu\nu}\partial_\rho\chi\ne0\quad\text{generically},\\
 [G_4\hbar]&=L^2,\quad G_{\rm proposed}=\zeta\frac{L_*^2}{\hbar}.
\end{align}
The final expression predicts $G$ only if $\zeta$ and $L_*$ are independently fixed, with no $G$ used in their definition.
\end{document}
